% likelihood_illustration.tex
\begin{tikzpicture}
\begin{axis}[
    width=7cm, height=5cm,
    xlabel={$y$},
    ylabel={$p(y \mid \vx;\, \pvec)$},
    xmin=-4, xmax=4,
    ymin=0, ymax=0.55,
    xtick=\empty,
    ytick=\empty,
    label style={font=\small},
    tick label style={font=\scriptsize},
    axis lines=left,
    clip=false,
]
 
% high likelihood Gaussian — centered on data
\addplot[thick, thwspetrol, solid, domain=-4:4, samples=50]
    {0.5*exp(-0.5*((x-0)/0.6)^2)/(0.6*sqrt(2*pi))};
 
% low likelihood Gaussian — shifted and wider
\addplot[thick, thwsorange, dashed, domain=-4:4, samples=50]
    {0.5*exp(-0.5*((x-2.5)/0.4)^2)/(0.4*sqrt(2*pi))};
 
% data points on x-axis
\addplot[only marks, mark=|, thwspetrol, mark size=4pt] coordinates {
    (-1.2, 0) (1.3, 0) (0.1, 0) (-0.5, 0) (0.9, 0) (-1.5, 0) (-0.3, 0) (-0.7, 0) (0.4, 0)
};
 
% annotations
\node[thwspetrol, font=\scriptsize, anchor=south] at (axis cs:0, 0.35)
    {high likelihood};
\node[thwsorange, font=\scriptsize, anchor=south] at (axis cs:2.5, 0.5)
    {low likelihood};
 
\end{axis}
\end{tikzpicture}